% Part: sets-functions-relations
% Chapter: relations
% Section: orders

\documentclass[../../../include/open-logic-section]{subfiles}

\begin{document}

\olfileid[pt]{sfr}{rel}{ord}
\olsection{Ordens}

\begin{explain}
Muitas das nossas comparações envolvem descrever certos objetos como
sendo ``menores que'', ``iguais a'' ou ``maiores que'' outros objetos,
sob algum aspecto. Isso envolve relações de \emph{ordem}. Mas há
vários tipos de relação de ordem. Por exemplo, algumas exigem que
quaisquer dois objetos sejam comparáveis, outras não. Algumas incluem a
identidade (como~$\le$) e outras a excluem (como~$<$). Ter aqui uma
taxonomia ajuda-nos.
\end{explain}

\begin{defn}[Pré-ordem]
Uma relação que é ao mesmo tempo reflexiva e transitiva chama-se uma
\emph{pré-ordem.}
\end{defn}

\begin{defn}[Ordem parcial]
Uma pré-ordem que também é antissimétrica chama-se uma
\emph{ordem parcial}.
\end{defn}

\begin{defn}[Ordem linear]\ollabel{def:linearorder}
Uma ordem parcial que também é conexa chama-se uma
\emph{ordem total} ou \emph{ordem linear.}
\end{defn}

\begin{ex}
Toda ordem linear é também uma ordem parcial, e toda ordem parcial é
também uma pré-ordem, mas as implicações recíprocas não valem. A relação
universal sobre~$A$ é uma pré-ordem, pois é reflexiva e transitiva. Mas,
se $A$ tem mais de um !!{element}, a relação universal não é
antissimétrica e, portanto, não é uma ordem parcial.
\end{ex}

\begin{ex}
Considere a relação \emph{não mais longa que} $\preccurlyeq$
sobre~$\Bin^*$: $x \preccurlyeq y$ sse $\len{x} \le \len{y}$. Trata-se de
uma pré-ordem (reflexiva e transitiva), e até conexa, mas não é uma
ordem parcial, pois não é antissimétrica. Por exemplo, $01
\preccurlyeq 10$ e $10 \preccurlyeq 01$, mas $01 \neq 10$.
\end{ex}

\begin{ex}
Uma ordem parcial importante é a relação $\subseteq$ sobre um conjunto de
conjuntos. Em geral não é uma ordem linear: se $a \neq b$ e
consideramos $\Pow{\{a, b\}} = \{\emptyset, \{a\}, \{b\}, \{a,b\}\}$,
vemos que $\{a\} \nsubseteq \{b\}$ e $\{a\} \neq \{b\}$ e $\{b\}
\nsubseteq \{a\}$.
\end{ex}

\begin{ex}
A relação de \emph{divisibilidade sem resto} nos dá uma
ordem parcial que não é linear. Para inteiros $n$ e~$m$, escrevemos
$n \mid m$ para significar que $n$ divide $m$ (exatamente), isto é, sse
existe algum inteiro~$k$ tal que $m = kn$. Sobre~$\Nat$, isso é uma ordem parcial,
mas não linear: por exemplo, $2 \nmid 3$ e também $3 \nmid
2$. Vista como relação sobre $\Int$, a divisibilidade é apenas uma
pré-ordem, pois não é antissimétrica: $1 \mid -1$ e $-1 \mid 1$
mas $1 \neq -1$.
\end{ex}

\begin{ex}
A relação de \emph{extensão} sobre um conjunto de sequências~$A^*$ é a
seguinte: $s \sqsubseteq s'$ sse $s = \emptyseq$ (a sequência
vazia), $s = s'$, ou $s = \tuple{s_1, \dots, s_n}$ e $s' =
\tuple{s_1, \dots, s_n, s_{n+1}, \dots, s_m}$. Se $s \sqsubseteq s'$
dizemos também que $s$ é um \emph{segmento inicial} de~$s'$. A
relação de extensão sobre $A^*$ é uma ordem parcial, mas não linear,
p.ex., se $a \neq b$, então $ab \not\sqsubseteq ba$ e $ba \not\sqsubseteq ab$.
\end{ex}

\begin{defn}[Ordem estrita]
Uma \emph{ordem estrita} é uma relação irreflexiva, assimétrica
e transitiva.
\end{defn}

\begin{defn}[Ordem linear estrita]\ollabel{def:strictlinearorder}
Uma ordem estrita que também é conexa chama-se
\emph{ordem total estrita} ou \emph{ordem linear estrita.}
\end{defn}

\begin{ex}
$\le$ é a ordem linear correspondente à ordem linear estrita~$<$.
$\subseteq$ é a ordem parcial correspondente à ordem estrita~$\subsetneq$.
\end{ex}

Qualquer ordem estrita $R$ sobre~$A$ pode transformar-se numa ordem parcial
adicionando a diagonal $\Id{A}$, isto é, adicionando todos os pares~$\tuple{x,
x}$. (Isto chama-se o \emph{fecho reflexivo} de~$R$.)
Reciprocamente, a partir de uma ordem parcial obtém-se uma ordem estrita
removendo~$\Id{A}$. Os dois resultados seguintes tornam isto preciso.

\begin{prop}\ollabel{prop:stricttopartial}
Se $R$ é uma ordem estrita sobre~$A$, então $R^+ = R \cup \Id{A}$ é uma
ordem parcial. Além disso, se $R$ é uma ordem linear estrita, então $R^+$ é
uma ordem linear.
\end{prop}

\begin{proof}
Suponhamos que $R$ é uma ordem estrita, isto é, $R \subseteq A^2$ e $R$ é
irreflexiva, assimétrica e transitiva. Seja $R^+ = R \cup \Id{A}$. Temos
de mostrar que $R^+$ é reflexiva, antissimétrica e transitiva.

$R^+$ é claramente reflexiva, pois $\tuple{x, x} \in \Id{A} \subseteq
R^+$ para todo $x \in A$.

Para mostrar que $R^+$ é antissimétrica, suponhamos, para redução ao absurdo, que $R^+xy$ e
$R^+yx$ mas $x \neq y$. Como $\tuple{x,y} \in R \cup \Id{A}$, mas
$\tuple{x, y} \notin \Id{A}$, temos de ter $\tuple{x, y} \in R$, isto é,
$Rxy$. Do mesmo modo,~$Ryx$. Mas isto contradiz a hipótese de que
$R$ é assimétrica.

Para estabelecer a transitividade, suponhamos que $R^+xy$ e $R^+yz$. Se tanto
$\tuple{x, y} \in R$ como $\tuple{y,z} \in R$, então $\tuple{x, z} \in
R$, pois $R$~é transitiva. Caso contrário, ou $\tuple{x, y} \in
\Id{A}$, isto é, $x = y$, ou $\tuple{y, z} \in \Id{A}$, isto é, $y = z$.
No primeiro caso, temos $R^+yz$ por hipótese, $x = y$, logo
$R^+xz$. Analogamente no segundo caso. Em qualquer caso, $R^+xz$; assim,
$R^+$ é também transitiva.

Quanto à cláusula ``além disso'', suponhamos que $R$ é também conexa.
Então, para todo $x \neq y$, ou $Rxy$ ou~$Ryx$, isto é, ou
$\tuple{x, y} \in R$ ou $\tuple{y, x} \in R$. Como $R \subseteq R^+$,
isto continua a valer para $R^+$, logo $R^+$ é também conexa.
\end{proof}

\begin{prop}\ollabel{prop:partialtostrict}
Se $R$ é uma ordem parcial sobre~$A$, então $R^- = R \setminus \Id{A}$ é uma
ordem estrita. Além disso, se $R$ é uma ordem linear, então $R^-$ é uma ordem
linear estrita.
\end{prop}

\begin{proof}
Deixamos como exercício.
\end{proof}

\begin{prob}
Apresente uma prova da \olref[sfr][rel][ord]{prop:partialtostrict}.
\end{prob}

O resultado simples a seguir estabelece que ordens lineares estritas
satisfazem uma propriedade no estilo da extensionalidade:

\begin{prop}\ollabel{prop:extensionality-strictlinearorders}
Se $<$ é uma ordem linear estrita sobre $A$, então:
\[
  (\forall a, b \in A)((\forall x \in A)(x < a \liff x < b) \lif a = b).
\]
\end{prop}

\begin{proof}
Suponhamos $(\forall x \in A)(x < a \liff x < b)$. Se $a < b$, então $a <
a$, o que contradiz o fato de $<$ ser irreflexiva; logo $a \nless b$.
De forma totalmente análoga, $b \nless a$. Logo $a = b$, pois $<$ é conexa.
\end{proof}

\end{document}
